\documentclass[10pt,a4paper]{amsart}
\usepackage{amsmath,amssymb,amsthm,graphicx}
% Retained preamble for the Stacks Project chapter morphisms.tex (Morphisms of Schemes).
% The chapter file pinned-original-source.tex carries only \input{preamble} on its first line, so no
% preamble text lives inside the chapter itself.  The Stacks Project uses one single shared preamble for
% every chapter; the snapshot stages it verbatim in sources/stacks-chow/preamble.tex.  The block below is
% copied mechanically, line for line, from that file's lines 14-18 (the Xy-pic setup) and lines 154-217
% (the native theorem declarations together with the mathematical macros).  Nothing was added, removed or
% redefined.  The wrapper class and the amsmath/amssymb/amsthm/graphicx base packages are supplied by the
% assembler, so the class/\externaldocument part of the shared preamble is not repeated here.
% For commutative diagrams we use Xy-pic
\usepackage[all]{xy}

% We use 2cell for 2-commutative diagrams.
\xyoption{2cell}
% Theorem environments.
%
\theoremstyle{plain}
\newtheorem{theorem}[subsection]{Theorem}
\newtheorem{proposition}[subsection]{Proposition}
\newtheorem{lemma}[subsection]{Lemma}

\theoremstyle{definition}
\newtheorem{definition}[subsection]{Definition}
\newtheorem{example}[subsection]{Example}
\newtheorem{exercise}[subsection]{Exercise}
\newtheorem{situation}[subsection]{Situation}

\theoremstyle{remark}
\newtheorem{remark}[subsection]{Remark}
\newtheorem{remarks}[subsection]{Remarks}

\numberwithin{equation}{subsection}

% Macros
%
\def\lim{\mathop{\mathrm{lim}}\nolimits}
\def\colim{\mathop{\mathrm{colim}}\nolimits}
\def\Spec{\mathop{\mathrm{Spec}}}
\def\Hom{\mathop{\mathrm{Hom}}\nolimits}
\def\Ext{\mathop{\mathrm{Ext}}\nolimits}
\def\SheafHom{\mathop{\mathcal{H}\!\mathit{om}}\nolimits}
\def\SheafExt{\mathop{\mathcal{E}\!\mathit{xt}}\nolimits}
\def\Sch{\mathit{Sch}}
\def\Mor{\mathop{\mathrm{Mor}}\nolimits}
\def\Ob{\mathop{\mathrm{Ob}}\nolimits}
\def\Sh{\mathop{\mathit{Sh}}\nolimits}
\def\NL{\mathop{N\!L}\nolimits}
\def\CH{\mathop{\mathrm{CH}}\nolimits}
\def\proetale{{pro\text{-}\acute{e}tale}}
\def\etale{{\acute{e}tale}}
\def\QCoh{\mathit{QCoh}}
\def\Ker{\mathop{\mathrm{Ker}}}
\def\Im{\mathop{\mathrm{Im}}}
\def\Coker{\mathop{\mathrm{Coker}}}
\def\Coim{\mathop{\mathrm{Coim}}}

% Boxtimes
%
\DeclareMathSymbol{\boxtimes}{\mathbin}{AMSa}{"02}

%
% Macros for moduli stacks/spaces
%
\def\QCohstack{\mathcal{QC}\!\mathit{oh}}
\def\Cohstack{\mathcal{C}\!\mathit{oh}}
\def\Spacesstack{\mathcal{S}\!\mathit{paces}}
\def\Quotfunctor{\mathrm{Quot}}
\def\Hilbfunctor{\mathrm{Hilb}}
\def\Curvesstack{\mathcal{C}\!\mathit{urves}}
\def\Polarizedstack{\mathcal{P}\!\mathit{olarized}}
\def\Complexesstack{\mathcal{C}\!\mathit{omplexes}}
% \Pic is the operator that assigns to X its picard group, usage \Pic(X)
% \Picardstack_{X/B} denotes the Picard stack of X over B
% \Picardfunctor_{X/B} denotes the Picard functor of X over B
\def\Pic{\mathop{\mathrm{Pic}}\nolimits}
\def\Picardstack{\mathcal{P}\!\mathit{ic}}
\def\Picardfunctor{\mathrm{Pic}}
\def\Deformationcategory{\mathcal{D}\!\mathit{ef}}

% --- appended 2026-10-12 (Chief): the project own definitions of the reference/slogan/history
% display environments, copied mechanically from the shared Stacks preamble lines 8-12
% (\usepackage{verbatim} plus the three \newenvironment lines, which the chapter preamble keeps
% verbatim from the project preamble). Nothing added, removed or redefined: these environments
% carry no mathematical content and the project itself typesets them as comments. Without these
% three lines a counted statement whose own text uses them fails with Environment slogan undefined.
\usepackage{verbatim}
\newenvironment{reference}{\comment}{\endcomment}
\newenvironment{slogan}{\comment}{\endcomment}
\newenvironment{history}{\comment}{\endcomment}
\usepackage{cleveref}
\crefname{originallemma}{original lemma}{original lemmas}
\Crefname{originallemma}{Original lemma}{Original lemmas}
\makeatletter
\newlabel{algebra-lemma-image-stable-specialization-closed}{{00HY}{source}{}{sourcepointer-1}{}}
\newlabel{algebra-lemma-image-stable-specialization-closed@cref}{{[originallemma][0][]00HY}{[1][0][]source}{}{}{}}
\makeatother
\crefname{originallemma}{original lemma}{original lemmas}
\Crefname{originallemma}{Original lemma}{Original lemmas}
\makeatletter
\newlabel{lemma-generalizations-lift-flat}{{03HV}{source}{}{sourcepointer-2}{}}
\newlabel{lemma-generalizations-lift-flat@cref}{{[originallemma][0][]03HV}{[1][0][]source}{}{}{}}
\makeatother
\crefname{originallemma}{original lemma}{original lemmas}
\Crefname{originallemma}{Original lemma}{Original lemmas}
\makeatletter
\newlabel{topology-lemma-lift-specializations-images}{{0065}{source}{}{sourcepointer-3}{}}
\newlabel{topology-lemma-lift-specializations-images@cref}{{[originallemma][0][]0065}{[1][0][]source}{}{}{}}
\makeatother
\crefname{originallemma}{original lemma}{original lemmas}
\Crefname{originallemma}{Original lemma}{Original lemmas}
\makeatletter
\newlabel{topology-lemma-open-closed-specialization}{{0062}{source}{}{sourcepointer-4}{}}
\newlabel{topology-lemma-open-closed-specialization@cref}{{[originallemma][0][]0062}{[1][0][]source}{}{}{}}
\makeatother
\title{A quasi-compact surjective flat morphism is submersive on the underlying topological spaces}
\author{The Stacks Project Authors}
\date{stacks-project:morphisms.tex:local-snapshot-20261012; GNU FDL 1.2 or later}
\begin{document}
\maketitle
\noindent\textit{Editorial note.} All delivered mathematical text is the author's own, copied from the pinned Stacks Project chapter file pinned-original-source.tex; no statement or proof step was rewritten, reordered or removed.\par\medskip
\noindent\textit{Editorial note.} Following the multi-claim convention for this source, source\_url is the source-level URL https://stacks.math.columbia.edu/ and the source label lemma-fpqc-quotient-topology is carried in claim\_key and here rather than in a per-claim tag URL.\par\medskip
\noindent\textit{Editorial note.} The preamble is the chapter's shared Stacks Project preamble, retained as sources/stacks-morphisms/retained-preamble.tex: a line-for-line excerpt of the shared Stacks preamble staged with this snapshot. No macro was invented or redefined, and this card adds no preamble text of its own.\par\medskip
\noindent\textit{Editorial note.} Statement span 4843-4853 and proof span 4855-4878 come from the halfA pool card operations/pool-cards/stacks-morphisms-halfA.json and were re-read in the pinned chapter file before this card was written.\par\medskip
\noindent\textit{Editorial note.} The counted statement contains a reference environment holding the author's two citations (SGA1 and EGA); the project's own shared preamble defines that environment as a comment environment, so no bibliography is attached and none is needed.\par\medskip
\noindent\textit{Editorial note.} The statement's 'resp. closed' clause and the proof's use of complements are the author's wording and are preserved.\par\medskip
\noindent\textit{Editorial note.} 4 cross-references occur inside the counted spans; they are listed in external\_references as navigation pointers only.\par\medskip
\noindent\textit{Allowed standard prerequisites (editorial).}\begin{itemize}
\item The question of openness and closedness of a subset of the target being local on the target, and quasi-compactness of the source over an affine base. Used at proof lines 3072-3079. Permits replacing the base by an affine open and covering the quasi-compact source by finitely many affine opens, including after replacing the source by their disjoint union.
\item Generalizations lift along flat morphisms, cited from the same chapter, and the topological lemma that images of sets stable under generalization are the complements of sets stable under specialization. Cited at proof lines 3090-3093. Shows that the complement of the candidate closed subset is stable under generalization, hence the subset itself is stable under specialization.
\item The algebra lemma that a subset of an affine scheme stable under specialization is closed. Cited at proof lines 3094-3095. Closes the affine case of the argument, from which the general case follows by the initial reductions.
\end{itemize}\medskip
\ifdefined\hypertarget\hypertarget{sourcepointer-1}{}\fi
\noindent\textit{Source reference.} Original lemma 00HY: focused-seeds/AGP/sources/stacks-snapshot-20261011/tags.txt, source lines 660--660. Same-chapter reference used in the retained author text at source line(s) 4877. It is supplied only as a navigation pointer so that the retained text compiles; the counted item delivers the author's own statement and the author's complete proof, and no part of the referenced item is copied in or substituted for the counted argument.\par\medskip
\ifdefined\hypertarget\hypertarget{sourcepointer-2}{}\fi
\noindent\textit{Source reference.} Original lemma 03HV: focused-seeds/AGP/sources/stacks-snapshot-20261011/tags.txt, source lines 4581--4581. Same-chapter reference used in the retained author text at source line(s) 1315, 4469, 4803, 4871, 14824. It is supplied only as a navigation pointer so that the retained text compiles; the counted item delivers the author's own statement and the author's complete proof, and no part of the referenced item is copied in or substituted for the counted argument.\par\medskip
\ifdefined\hypertarget\hypertarget{sourcepointer-3}{}\fi
\noindent\textit{Source reference.} Original lemma 0065: focused-seeds/AGP/sources/stacks-snapshot-20261011/tags.txt, source lines 233--233. Same-chapter reference used in the retained author text at source line(s) 4872. It is supplied only as a navigation pointer so that the retained text compiles; the counted item delivers the author's own statement and the author's complete proof, and no part of the referenced item is copied in or substituted for the counted argument.\par\medskip
\ifdefined\hypertarget\hypertarget{sourcepointer-4}{}\fi
\noindent\textit{Source reference.} Original lemma 0062: focused-seeds/AGP/sources/stacks-snapshot-20261011/tags.txt, source lines 230--230. Same-chapter reference used in the retained author text at source line(s) 4875. It is supplied only as a navigation pointer so that the retained text compiles; the counted item delivers the author's own statement and the author's complete proof, and no part of the referenced item is copied in or substituted for the counted argument.\par\medskip
\begin{lemma}
\label{lemma-fpqc-quotient-topology}
\begin{reference}
\cite[Expose VIII, Corollaire 4.3]{SGA1} and
\cite[IV, Corollaire 2.3.12]{EGA}
\end{reference}
Let $f : X \to Y$ be a quasi-compact, surjective, flat morphism.
A subset $T \subset Y$ is open (resp.\ closed) if and only
$f^{-1}(T)$ is open (resp.\ closed). In other words, $f$ is
a submersive morphism.
\end{lemma}

\begin{proof}
The question is local on $Y$, hence we may assume that $Y$ is affine.
In this case $X$ is quasi-compact as $f$ is quasi-compact.
Write $X = X_1 \cup \ldots \cup X_n$ as a finite union of affine opens.
Then $f' : X' = X_1 \amalg \ldots \amalg X_n \to Y$ is a surjective
flat morphism of affine schemes. Note that for $T \subset Y$ we have
$(f')^{-1}(T) = f^{-1}(T) \cap X_1 \amalg \ldots \amalg f^{-1}(T) \cap X_n$.
Hence, $f^{-1}(T)$ is open if and only if $(f')^{-1}(T)$ is open.
Thus we may assume both $X$ and $Y$ are affine.

\medskip\noindent
Let $f : \Spec(B) \to \Spec(A)$ be a surjective morphism of affine schemes
corresponding to a flat ring map $A \to B$. Suppose that $f^{-1}(T)$ is
closed, say $f^{-1}(T) = V(J)$ for $J \subset B$ an ideal. Then
$T = f(f^{-1}(T)) = f(V(J))$ is the image of $\Spec(B/J) \to \Spec(A)$
(here we use that $f$ is surjective). On the other hand, generalizations
lift along $f$ (Lemma \ref{lemma-generalizations-lift-flat}).
Hence by Topology, Lemma \ref{topology-lemma-lift-specializations-images}
we see that $Y \setminus T = f(X \setminus f^{-1}(T))$ is stable under
generalization. Hence $T$ is stable under specialization
(Topology, Lemma \ref{topology-lemma-open-closed-specialization}).
Thus $T$ is closed by
Algebra, Lemma \ref{algebra-lemma-image-stable-specialization-closed}.
\end{proof}

\end{document}
